\subsection{A finite sparse characteristic effect that disappears later}

\begin{proposition}[Proved: targetwise separation need not persist]
\label{C4:sparse-local-global}
In $\Cdelta$, let $V$ be the singleton and let $Q$ have
ordered vertices $a<b<c<x<y$ and edges $\{a,x,y\}$ and
$\{b,c,y\}$. Consider the three strong arrows $V\to Q$
selecting $a,b,c$. There is a strong target $R$ with
$24$ vertices and $14$ edges such that this family has an
affine equalizer over $\mathbb F_2$ in $R$, but none over
$\Q$ in $R$. Nevertheless the same family has an integral
affine equalizer at a larger target in $\Cdelta$.
\end{proposition}
\begin{proof}
Use ten selected vertices in the order
\[
 B<A<C<E<z<D<F<J<H<I
\]
and the following seven increasing triples:
\[
 (z,D,H),\ (A,z,I),\ (B,E,z),\ (A,E,J),\
 (C,D,J),\ (C,F,H),\ (B,F,I).
 \tag{SL.1}
\]
For each triple $(a_i,b_i,c_i)$, add private vertices
$x_i<y_i$, put all private vertices above the ten selected
vertices, and add the edges $a_ix_iy_i,b_ic_iy_i$.
Distinct triples meet in at most one vertex.

Here is a complete finite predimension certificate. For a
subset $S$ of the ten selected vertices, minimizing over
the private vertices of one gadget gives contribution
$-2$ if all three selected vertices of that gadget belong
to $S$, contribution $-1$ if $b_i,c_i\in S$ but $a_i\notin S$,
and contribution zero otherwise. Thus the minimum is
\[
 m(S)=2|S|-2\#\{i:\{a_i,b_i,c_i\}\subseteq S\}
       -\#\{i:\{b_i,c_i\}\subseteq S,\ a_i\notin S\}.
 \tag{SL.2}
\]
The minima of (SL.2) at successive cardinalities
$|S|=0,1,\ldots,10$ are
\[
 0,2,3,4,5,5,6,6,6,6,6.
 \tag{SL.3}
\]
The exact certificate \texttt{code/C4\_sparse\_local\_exception.py}
checks all $1024$ selected subsets and all four private
choices in each gadget; therefore it covers all $2^{24}$
actual subsets. A nonempty subset containing no selected
vertex has no edge and positive predimension. Equations
(SL.2)--(SL.3) prove positivity. Each proposed copy of $Q$
has predimension four. Adding a selected vertex gives
predimension at least five; with its selected triple fixed,
adding other private vertices also strictly raises
predimension, since every other gadget meets the triple
in at most one point. Consequently all seven copies are
strong, and $\delta(R)=6$.

They are all the induced copies of this ordered $Q$ in $R$.
Its two edges must meet in their greatest vertex. Two
edges from different gadgets can meet only at a selected
vertex, which is below their private vertices. They cannot
form such a copy. Within a gadget the increasing order
forces the displayed identification with $Q$.

Let $D$ have one column per copy, recording
$(b_i-a_i,c_i-a_i)$ in two copies of the free module on
the selected vertices. Every vertex other than $z$ occurs
twice in (SL.1), always in the same position; $z$ occurs
once in each position. Hence the sum of the seven columns
is zero modulo two, with augmentation $7=1$ modulo two.
Over $\Q$, a vector $(t_i)_{i=0}^6$ in the kernel must
satisfy $t_i+t_j=0$ whenever two triples share a vertex
other than $z$. The resulting graph of the seven triples
is connected and contains the triangle $0,4,5,0$.
It follows that every $t_i=0$. The exact integer matrix
has Smith factors $1,1,1,1,1,1,2$, as independently
recorded by the certificate. There is therefore no
rational affine equalizer at this target.

We next prove that the targetwise obstruction disappears.
Let $M$ be the strong Fraisse limit and let $\Omega$ be
its strong-singleton orbit. Form the abelian group with
generators $u_t,v_t$ for $t\in\Omega$ and one generator
$h$, with the relations
\[
 u_b-u_a+v_c-v_a=h
 \quad\text{for every strong ordered copy }(a,b,c,x,y)
       \text{ of }Q\text{ in }M.
 \tag{SL.4}
\]
Write $\psi_t=u_t+v_t$. A three-edge star with edges
$axy,a'x'y,bcy$, ordered so that both $a,a'$ are below
$b<c$ and both $x,x'$ lie between $c$ and $y$, has
seven vertices and predimension five. Its two displayed
copies of $Q$ are strong. Subtracting their relations
gives $\psi_a=\psi_{a'}$.

The pair $a,a'$ is not strong in this star; this is why
one star does not suffice to compare an arbitrary strong
pair. Glue two such stars along an outer minimum, obtaining
outer minima $a_0,a_1,a_2$. All other vertices may be put
above these three minima in an order satisfying both
stars. Their free amalgam is positive, and its four
specified copies of $Q$ are strong. Its endpoint pair
$\{a_0,a_2\}$ is strong: if a subset containing this pair
omits $a_1$, its intersections with the two stars each
have predimension at least two, strictly larger when
enlarged beyond the endpoint; if it contains $a_1$, each
intersection contains the two outer minima and has
predimension at least four, so their union has
predimension at least $4+4-2=6>4$. This also has an
independent exhaustive certificate on all $8191$ nonempty
subsets of the thirteen-vertex bridge.

The extension property now gives $\psi_s=\psi_t$ for any
strong pair $s,t$ in $M$. For an arbitrary pair in
$\Omega$, add one singleton freely over a finite strong
set containing it. Each old point and the new point form
a strong pair, so $\psi$ is constant on all of $\Omega$.
Every strong edge $b<c<y$ extends to a strong copy of $Q$
by adding $a<b$ and $c<x<y$. Substituting the constant
$\psi$ into (SL.4) gives
\[
 u_b-u_c=h\qquad\text{for every strong ordered edge }b<c<y.
 \tag{SL.5}
\]
The integral pair affine-equalizer property of the total
category $\Cdelta$, applied to the first two singleton
arrows into an edge, makes a finite integer combination
of (SL.5) have left side zero and right side $h$.
Therefore $h=0$ in the group presented by (SL.4).

Membership in the subgroup generated by the relations
uses only finitely many relations. Writing $h$ as their
integer combination gives a finite integer combination
of copies of $Q$ whose two marginal differences vanish
and whose augmentation is $1$, after changing all signs
if necessary. A finite strong set containing their images
is a target for an integral affine equalizer. Taking a
common successor with $R$ if desired makes it a target
extending the original finite-stage certificate.
\end{proof}

\paragraph{Scope.}
This example separates a matrix calculation at one target
from the existence of an affine equalizer somewhere in the
category. It neither gives a familywise modular-versus-rational
separation in a total category nor changes the status of
the total restriction of C.4.
